% Integración sucesora 2026-09-26: el lector longitudinal precede a G.
% La realización circular anterior se conserva como corolario, no como generador.
\section{Action, longueur structurelle, gravité et information thermique}
\label{v:ant:sec:gravity}

\subsection{Retour enregistré et construction de la longueur}

Nous écrivons $\alpha_{\HMT}$ pour la coordonnée dodécaphasique déjà engendrée
---notée $\alpha$ dans les sections précédentes---. Aucune valeur
métrologique de $G$ ni aucune longueur de Planck tabulée n’intervient ici. Le
registre marqué réunit
\begin{equation}
 N=12\cdot4\cdot120=5760
 \label{v:eq:cardinal-longitud-alpha}
\end{equation}
positions. Soient $u_\Omega=N^{-1/2}\mathbf1_\Omega$,
$Q_\Omega=I-u_\Omega u_\Omega^*$ et
$\mathsf B:V_U\to V_\Omega$ le retour de Hadamard typé, avec
\begin{equation}
 \mathsf B^*\mathsf B=\mathsf B\mathsf B^*=\frac14I,
 \qquad \mathsf C=Q_\Omega\mathsf B=\mathsf BQ_U.
 \label{v:eq:retorno-hadamard-longitud}
\end{equation}
Pour le pincement diagonal
$\Delta_\Omega(A)=\sum_{i\in\Omega}E_iAE_i$, on obtient
\begin{equation}
 \Delta_\Omega(\mathsf C\mathsf C^*)
 =\frac{N-1}{4N}I=r_NI,
 \qquad
 r_N=\frac{5759}{23040},
 \label{v:eq:residuo-longitud-alpha}
\end{equation}
tandis que le canal complémentaire conserve $1/(4N)$. En effet,
$\mathsf C\mathsf C^*=Q_\Omega/4$ et chaque entrée diagonale de
$Q_\Omega$ vaut $1-1/N$. Le pincement détermine la réponse
longitudinale ; l’inscription élargie et son complément conservent la
mémoire que la lecture diagonale ne distingue pas.

Soit $\mathcal L_L$ la droite orientée de longueur et soit $L_*>0$ sa
section de référence transportée par la cochaîne longitudinale. Cette
section n’est pas identifiée tacitement à la base native d’arête :
sa relation avec l’horloge est fixée ensuite par la réalisation circulaire.
La norme du conoyau $H_4$ déjà construite apporte
$\mathrm N_{H_4}(\alpha_{\HMT})=\alpha_{\HMT}^{16}$. La composition
linéaire du retour enregistré détermine donc
\begin{equation}
 \boxed{
 L_\alpha=\alpha_{\HMT}^{16}r_NL_*
 =\alpha_{\HMT}^{16}\frac{5759}{23040}L_* .}
 \label{v:eq:longitud-alpha}
\end{equation}
La puissance, le coefficient d’incidence et la section dimensionnelle ont
des origines séparées et explicites. $L_\alpha$ est une sortie longitudinale
antérieure à la lecture gravitationnelle.

\subsection{Composition opératorielle et détermination unique de la gravité}

Soit $\mathcal H\ne\{0\}$ une fibre de Hilbert finie et soit $Y=UY_0U^*>0$
le caractère positif et inversible d’une route, avec son transport, son feuillet,
son orientation et sa mémoire conservés. Pour une section d’action
$\hbar_a>0$, $a\in\{\mathrm{pre},\mathrm{ret}\}$, et la vitesse
$c=c_{\mathrm{int}}=c_{\mathrm{vac}}=\widehat c\,u_C>0$ déjà construite dans
\eqref{v:eq:identidad-velocidad}, posons
\begin{equation}
 Q_{L,a}=\frac{\hbar_ac}{L_\alpha},\qquad
 H_a=Q_{L,a}Y,\qquad M_a=\frac{H_a}{c^2},\qquad
 \Lambda_a=L_\alpha Y^{-1},\qquad R_a=L_\alpha Y.
 \label{v:eq:lectores-longitud-accion}
\end{equation}
Le secteur nul est maintenu séparé et n’est pas inversé.

\begin{theorem}[Unicité du coefficient gravitationnel dans la réalisation radiale]
\label{v:thm:gravedad-unica-alpha}
Les lecteurs précédents satisfont
\begin{equation}
 H_a\Lambda_a=\hbar_ac\,I,\qquad
 R_a\Lambda_a=L_\alpha^2I.
 \label{v:eq:dos-relaciones-gravedad}
\end{equation}
Si la lecture physique déclarée du rayon réduit est
$R_a=(G_a/c^2)M_a$, alors il existe un unique coefficient compatible et c’est
\begin{equation}
 \boxed{G_a=\frac{c^3L_\alpha^2}{\hbar_a}},
 \qquad
 \boxed{\frac{G_a\hbar_a}{c^3}=L_\alpha^2}.
 \label{v:eq:gravedad-unica-alpha}
\end{equation}
\end{theorem}
\begin{proof}
Les définitions donnent
$H_a\Lambda_a=Q_{L,a}L_\alpha I=\hbar_acI$ et
$R_a\Lambda_a=L_\alpha^2I$. En éliminant $\Lambda_a$,
\[
 R_a=\frac{L_\alpha^2}{\hbar_ac}H_a.
\]
D’autre part,
$(G_a/c^2)M_a=G_aH_a/c^4$. Comme $H_a$ est inversible sur la fibre
non nulle, les deux expressions coïncident si et seulement si
$G_a=c^3L_\alpha^2/\hbar_a$. La même simplification montre que tout
autre $G_a'$ qui conserve les deux relations satisfait $G_a'=G_a$.
\end{proof}

La preuve utilise conjointement les deux lecteurs sur le même caractère ; elle
n’identifie pas la compression de $Y^{-1}$ à l’inverse d’une compression. Une
variation de route ou une élimination de mémoire qui agit de la même manière
sur $H_a$ et $R_a$ conserve leur proportionnalité et n’introduit aucun nouveau
facteur scalaire. L’analyse dimensionnelle confirme les exposants
$c^3L_\alpha^2\hbar_a^{-1}$, mais c’est
\eqref{v:eq:dos-relaciones-gravedad}, et non ce contrôle dimensionnel, qui
fixe le coefficient adimensionnel à un.

\subsection{L’écriture circulaire comme corollaire}

Le calendrier détermine $T_{108}=108t_0$ et
$\omega_{108}=2\pi_{\HMT}/T_{108}$. Dans la réalisation circulaire du même
lecteur longitudinal,
\begin{equation}
 \omega_{108}=\frac{c}{L_\alpha},\qquad
 R_{108}=\frac{c}{\omega_{108}}
 =\frac{54}{\pi_{\HMT}}ct_0=L_\alpha,
 \qquad
 t_0=\frac{\pi_{\HMT}L_\alpha}{54c}.
 \label{v:eq:compatibilidad-reloj-longitud}
\end{equation}
L’égalité déclare la compatibilité entre l’avance relevée de l’horloge
et la section longitudinale ; elle n’utilise pas $G_a$ pour choisir $L_\alpha$.
Sa substitution dans \eqref{v:eq:gravedad-unica-alpha} produit
\begin{equation}
 G_a=\left(\frac{54}{\pi_{\HMT}}\right)^2
 \frac{c^5t_0^2}{\hbar_a}.
 \label{v:eq:gravedad-circular-corolario}
\end{equation}
Par conséquent, dans l’expression paramétrée
$G_a(\vartheta)=\vartheta c^5t_0^2/\hbar_a$, la valeur
$\vartheta=(54/\pi_{\HMT})^2$ est un corollaire de la longueur construite,
et non un sélecteur circulaire indépendant.

Avec
$E_{108,a}=E_{P,a}=Q_{L,a}=\hbar_a\omega_{108}$ et
$m_{108,a}=E_{108,a}/c^2$ sont également conservées
\[
 \frac{\hbar_a}{m_{108,a}c}=R_{108}=L_\alpha,
 \qquad
 \frac{h_a}{m_{108,a}c}=2\pi_{\HMT}L_\alpha.
\]
Enfin, les deux orientations d’action satisfont
\begin{equation}
 \frac{G_{\mathrm{pre}}}{G_{\mathrm{ret}}}
 =\frac{\hbar_{\mathrm{ret}}}{\hbar_{\mathrm{pre}}}
 =R_{\mathrm{act}}^{-1},
 \qquad
 \ell_P^2:=L_\alpha^2
 =\frac{G_a\hbar_a}{c^3}.
 \label{v:eq:reciprocidad-gravedad-accion}
\end{equation}
L’unité surfacique procède ainsi de $L_\alpha^2$ ; la réciprocité ultérieure
entre action et gravité conserve cette aire lors du changement d’orientation.
